% 7
\begin{frame}{Reinforcement learning, in one slide}
\begin{columns}[T,onlytextwidth]
\column{0.54\textwidth}
\begin{center}
\begin{tikzpicture}[>=Latex,thick]
  \tikzset{node/.style={draw=DeepNavy,fill=SoftGray!50,rounded corners,minimum width=2cm,minimum height=.7cm,align=center}}
  \node[node] (s) at (0,1.6) {state $s_t$};
  \node[node] (a) at (3.0,1.6) {action $a_t$};
  \node[node] (env) at (3.0,0) {environment};
  \node[node] (r) at (0,0) {reward $r_t$};
  \draw[->] (s) -- node[above]{policy $\pi$} (a);
  \draw[->] (a) -- (env);
  \draw[->] (env) -- node[below]{observation} (r);
  \draw[->] (r) -- (s);
\end{tikzpicture}
\end{center}
\column{0.42\textwidth}
\small
\begin{itemize}
  \item \term{Policy}: how the agent chooses actions.
  \item \term{Rollout}: a sequence of states, actions, rewards.
  \item \term{Credit assignment}: which action caused later success?
  \item \term{Verifiable reward}: reward that can be computed by tests or checks.
\end{itemize}
\end{columns}

\vspace{0.45em}
\begin{block}{Coding is unusually attractive for RL}
A patch can be graded by executing real tests. That gives stronger feedback than chat preference alone.
\end{block}
\end{frame}

% 8
\begin{frame}{Model-free, model-based, and world-model RL}
\small
\begin{tabular}{p{0.24\textwidth}p{0.31\textwidth}p{0.35\textwidth}}
\toprule
\textbf{Approach} & \textbf{How it learns} & \textbf{Cost / failure mode} \\
\midrule
Model-free RL & Learn policy/value from real or logged rollouts. & Needs many interactions; reward can be sparse. \\
Model-based RL & Learn a dynamics model $\hat{P}(s'|s,a)$ and plan inside it. & Rollouts are cheap; model error compounds. \\
World models & Learn compressed latent dynamics from rich observations. & Great for imagination; must choose what information to keep. \\
Executable forks & Do not learn the transition for code; run it in forked worlds. & Real compute/storage; boundary must be faithful. \\
\bottomrule
\end{tabular}

\vspace{0.8em}
\centering\large The question is not \emph{model or no model}; it is \hilite{which future is cheaper to trust?}
\end{frame}

% 9
\begin{frame}{There are two ways to make futures cheap}
\begin{columns}[T,onlytextwidth]
\column{0.47\textwidth}
\begin{block}{Learn a world}
\centering
\large Dreamer / Contrastive World Models
\end{block}
\normalsize
\begin{center}
$o_t \rightarrow z_t \rightarrow \widehat{s}_{t+1},\widehat{s}_{t+2},\ldots$
\end{center}
\begin{itemize}
  \item compact learned latent state
  \item imagine many futures cheaply
  \item useful when interaction is slow, unsafe, or scarce
  \item pay approximation / learned-dynamics error
\end{itemize}
\small CWM's twist: retain future-predictive information, avoid reconstructing irrelevant visual detail.

\column{0.47\textwidth}
\begin{block}{Fork the world}
\centering
\large executable software environment
\end{block}
\normalsize
\begin{center}
$S_0 \rightarrow \{S_0^1,S_0^2,\ldots,S_0^N\}$
\end{center}
\begin{itemize}
  \item explicit executable state
  \item branch real repo / process futures
  \item useful when reset is cheap
  \item pay compute / storage, not learned transition-model error
\end{itemize}
\small Exact only inside the state and capability boundary we capture.
\end{columns}
\source{Dreamer 4: Hafner et al., arXiv:2509.24527. Contrastive World Models: Li, arXiv:2609.22175.}
\end{frame}

% 10
\begin{frame}{For code, the best world model is often the world}
\centering
\begin{tikzpicture}[>=Latex,thick]
  \tikzset{box/.style={draw=DeepNavy,fill=SoftGray!45,rounded corners,minimum width=2.2cm,minimum height=.82cm,align=center}}
  \tikzset{child/.style={draw=SignalCyan,fill=SignalCyan!8,rounded corners,minimum width=2.0cm,minimum height=.68cm,align=center}}
  \node[box] (s0) at (-5.2,0) {$S_0$\\snapshot};
  \node[child] (c1) at (-2.1,1.45) {patch A};
  \node[child] (c2) at (-2.1,.48) {patch B};
  \node[child] (c3) at (-2.1,-.49) {patch C};
  \node[child] (cn) at (-2.1,-1.46) {patch N};
  \node[box] (run) at (1.15,0) {execute\\real tests};
  \node[box] (red) at (4.05,0) {reduce\\evidence};
  \node[box,fill=WarmAmber!20] (prom) at (6.75,0) {promote\\once};
  \draw[->] (s0) -- node[above,sloped]{fork} (c1.west);
  \draw[->] (s0) -- (c2.west); \draw[->] (s0) -- (c3.west); \draw[->] (s0) -- (cn.west);
  \draw[->] (c1.east) -- (run.west); \draw[->] (c2.east) -- (run.west); \draw[->] (c3.east) -- (run.west); \draw[->] (cn.east) -- (run.west);
  \draw[->] (run) -- (red); \draw[->] (red) -- (prom);
  \node[color=Muted] at (-2.1,-2.28) {same past, different interventions};
  \node[color=Muted] at (6.75,-1.25) {burn children};
\end{tikzpicture}
\vspace{0.35em}
\Large These are \hilite{executable counterfactuals}: branch from the same past, intervene differently, observe the actual software transition.
\source{Related systems: Firecracker (NSDI 2020), DeltaBox (arXiv:2605.22781), Shepherd (arXiv:2605.10913).}
\end{frame}

% 11
\begin{frame}{Snapshot versus latent: the sufficiency question}
\small
\begin{tabular}{p{0.28\textwidth}p{0.30\textwidth}p{0.30\textwidth}}
\toprule
 & \textbf{World-model latent $z$} & \textbf{Executable snapshot $S$} \\
\midrule
Goal & Predict useful futures & Continue execution faithfully \\
Size & Compact, learned & Explicit, often overcomplete \\
Loss & Deliberately lossy & Lossless-ish inside a boundary \\
Error & Transition-model error & Boundary / reset / nondeterminism error \\
Use & Imagine rollouts & Execute branches \\
\bottomrule
\end{tabular}

\vspace{0.8em}
\begin{block}{Common question}
What information must survive so the future remains meaningful?
\end{block}
\centering
\large A snapshot is not a learned latent. It is an executable sufficient state.
\end{frame}
